% Meta-monografia de exemplo genérico de uso da classe delaetex.cls
% Copyright (C) 2004..2016 Walter Fetter Lages <fetter@ece.ufrgs.br>
%
% This file was adapted from:
% Meta-monografia de exemplo genérico de uso da classe deletex.cls
% Copyright (C) 2004 Walter Fetter Lages <w.fetter@ieee.org>
%
% This is free software, distributed under the GNU GPL; please take
% a look in `deletex.cls' to see complete information on using, copying
% and redistributing these files
%
%\documentclass[repeatfields,openright,overleaf,nomicrotype]{tcc}
\documentclass[repeatfields,xlists,xpacks,oneside,yearsonly]{ufrgscca}

%%%%%%%%%%%%%%%%%%%%%%
%%%%%%%%%%%%%%%%%%%%%%
%%
%% TCC Class options
%%
%% english   => paper/report in english
%%     \documentclass[english]{tcc}
%%
%% repeatfields => (abnt) references with same authors have their names repeated.
%%												otherwise only the first reference of them will be printed.
%%     \documentclass[repeatfields]{tcc}
%%
%% openright => force the start of a chapter at odd pages.
%%     \documentclass[openright]{tcc}
%%
%% oneside => single side document
%%     \documentclass[oneside]{tcc}
%%
%% NOTE : openright, oneside :: IF you are going to print the doc, single sided, the best is to just use
%%                              the [oneside] option.
%%                              IF you are going to print the doc, double sided, the best is to just use the
%%                              the [openright] option.
%%
%% relnum => figures/tables counters are reset with each chapter.
%%     \documentclass[relnum]{tcc}
%%
%% showframe => (for drafts only) auxiliary visual guide to page frame.
%%     \documentclass[showframe]{tcc}
%%
%% showlabels =>  (for drafts only) prints labes/citeref strings on the side bar.
%%     \documentclass[showlabels]{tcc}
%%
%% report => for simple report generation (only first cover page, no description or approval pages).
%%     \documentclass[report]{tcc}
%%
%% internship => fields adjustment for ´internship report´
%%     \documentclass[report,internship]{tcc}
%%
%% NOTE: with the report option,
%%                       the \course{} name is optional (one can remove it with \courseundef )
%%						 the \authorinfo{}{} is mandatory
%%       without the report option,
%%                       the \course{} name is mandatory
%%						 the \authorinfo{}{} is optional for the main work,
%%											but mandatory for many forms
%%
%%
%%%%%%%%%%%%%%%%%%%%%%
%%%%%%%%%%%%%%%%%%%%%%

%
% inicio do documento
%
%\RenewDocumentCommand{\fullcite}{}{}
%\RenewDocumentCommand{\cite}{}{}
%\let\MakeUppercase\uppercase
%\RenewDocumentCommand{\MakeUppercase}{m}{\uppercase{#1}}

% Pacotes
\usepackage{biblatex}
\usepackage{xcolor}
\usepackage{subfiles}
\usepackage{amsmath}
\usepackage{cancel}

% Bib
\addbibresource{referencias.bib}



% Cores especiais %
\definecolor{anotacao}{RGB}{0, 160, 50}
\definecolor{corrigir}{RGB}{210, 0, 0}
\definecolor{estranho}{RGB}{210, 70, 0}
\definecolor{att}{RGB}{0, 0, 230}

% Para escrever no escuro %
\makeatletter
\newcommand{\darkmode}{%
	\AtBeginDocument{
		\definecolor{dark-bk}{RGB}{31, 31 31}
		\color{white}\global\let\default@color\current@color
		\pagecolor{dark-bk}
	}
}
\makeatother


% Controlador  tema escuro %
%\darkmode

\begin{document}
% IMPORTAR DOCUMENTOS INICIAIS
\maketitle
\notoc\chapter{Agradecimentos}

Aos meus pais, Patrício e Eliane, e minha irmã Nicolle pelo apoio incondicional e sustentação emocional proporcionados durante todo o processo de graduação, mesmo apesar da distância, amo vocês.

A toda minha família incluindo tios, primos, avós, e em especial Luciano e Bruno, que também passaram pela experiência de se formar no estado gaúcho e foram grande fonte de companhia ao longo dos anos.

A todos os meus amigos e companheiros de curso, que formaram um verdadeiro grupo de apoio, proporcionando diversão e alívio em momentos difíceis.

Ao meu orientador e professor Mário pela alta preocupação e dedicação ao auxílio neste trabalho e na minha formação.

A todos os meus colegas de bolsa de pesquisa e de estágio, que proporcionaram um conhecimento externo e adicional ao ensinado pela universidade.

Essa conquista não seria possível sem vocês, muito obrigado. 

\begin{abstract}

No âmbito deste trabalho, busca-se o desenvolvimento de um sistema embarcado destinado a quadricópteros, visando controlar sua dinâmica com base em trajetórias predefinidas para operações autônomas de entrega. O sistema é modelado, empregando-se malhas de controle digital para gerenciar os estados que influenciam os motores da aeronave não tripulada (UAV, \textit{unmanned aerial vehicle}), utilizando o \textit{feedback} proveniente de uma unidade de medição inercial (IMU, \textit{inertial measurement unit}), um sistema de posicionamento global (GPS, \textit{global positioning system}) e demais sensores. Para determinar a configuração ideal do controlador para essa finalidade, são avaliadas duas abordagens distintas, o controlador PID e a linearização via realimentação de estados. Também é desenvolvido um algoritmo de voo que alimenta os controladores calculados para realização de entregas autônomas e posterior comparação. Ambas as técnicas são implementadas em um ambiente de simulação desenvolvido pelo autor, baseado na linguagem C++. Os resultados mostram que ambas as técnicas obtiveram bom comportamento para a aplicação vigente tanto para o algoritmo de controle quanto para o algoritmo de voo.

\end{abstract}
\begin{otherabstract}{Project, UAV, Autonomous Delivery, Automation and Control, Flight Dynamics, Quadcopters, Feedback linearization}

Within the scope of this work, the development of an embedded system intended for quadcopters is pursued, aiming to control their dynamics based on predefined trajectories for autonomous delivery operations. The system is modeled, employing digital control loops to manage the states influencing the motors of the unmanned aerial vehicle (UAV), utilizing feedback from an inertial measurement unit (IMU) and a global positioning system (GPS). To determine the optimal controller configuration for this purpose, two distinct approaches are evaluated: the PID controller and state feedback linearization. Additionally, a flight algorithm is developed to supply the computed controllers for performing autonomous deliveries and subsequent comparison. Both techniques are implemented in a simulation environment developed by the author, based on the C++ language. The results show that both techniques achieved good performance for the current application in both the control algorithm and the flight algorithm.
\end{otherabstract}
\setcounter{tocdepth}{3}
\listoffigures
\listoftables
\begin{listofabbrv}{PPGEE}
\item[UAV] Unmanned Aerial Vehicle (Veículo Aéreo Não-tripulado)
\item[IMU] Inertial Measurement Unit (Unidade de Mensuração Inercial)
\item[GPS] Global Positioning System (Sistema de Posicionamento Global)
\item[VTOL] Vertical Take-Off and Landing (Decolagem e Aterrisagem Verticais)
\item[PID] Controlador proporcional, integral e derivativo
\item[LRE] Linearização por Realimentação de Estados
\end{listofabbrv}
\tableofcontents

% IMPORTAR DOCUMENTOS PRINCIPAIS
\subfile{capitulos/introducao/}
\subfile{capitulos/descricao/}
\subfile{capitulos/modelagem/}
\subfile{capitulos/controle_atitude/}
\subfile{capitulos/algoritmo_voo/}
\subfile{capitulos/analise/}
\subfile{capitulos/conclusao/}


\printbibliography

\appendix
\subfile{capitulos/ambiente_simulacao/}


\end{document}

